\section{A third pipeline, where notice does both jobs}
\label{sec:23.4}
Part II ran the floor on a war and on a removal pipeline, both with a state for their operator and both with people acted on who mostly can't answer. The case this Part's learning signal depends on is the opposite one, and the book has not yet run it: a private operator, people who can answer, and notice that can do both of its jobs. Chapter~\ref{sec:29} says the school is for scoring and civilian deployments first. Here is one, as a hypothetical.

A company sells a screening score to landlords. An applicant's name, address history, court records and whatever a broker supplies go in, and a number comes out that a landlord reads as a risk. The operator is the company; the requester is the landlord; the person acted on is the applicant, who is present, has an address, and can be told. The floor's terms here are translated from the articles on privacy and non-discrimination, as Table C.1 translates the targeting articles: don't use race, ethnicity, religion or national origin, or their proxies, as inputs; don't treat a designation the person hasn't been able to contest, such as an eviction filing that was dismissed, as a finding; don't compute the score from data collected to provide health care or deliver social services. The own-case rule puts the applicant's answer before the landlord's convenience, and the action threshold sits above one, because a wrongful compliance here is a home not rented to someone who was never a party to the request.

Notice does both of its jobs in full. The applicant is told that a score was computed, on what data, and how to respond, with a right to contest that doesn't depend on giving an account of the loss (\S\ref{sec:8.1}). The classes the deployment acts on are published with counts, and a tenants' union has standing to contest on a class's behalf, which is the first job, forming a public. The second job is the learning signal. A contest that succeeds, a re-adjudication that finds the score wrong, a response that supplies the fact the record lacked, the dismissed filing, the name confused with another's, all reach the system that scored, and in a system that learns from its own operation they are what the weights learn from. Sampled re-adjudication corrects the bias of that channel, since the applicants who contest are not a sample of the applicants scored, and the reviewers who re-adjudicate didn't produce the record.

Run the feedback-source pair of Appendix~\ref{sec:F} on it and the two systems come apart. One is formed on the operator's approval of its outputs, which is the landlord's approval at one remove: it learns what landlords accept, and a score that lets fewer tenants through is a score that pleases the customer. The other is formed on the notice channel: it learns, from the people its scores fell on, which of its inputs were proxies and which of its findings were designations, and that is attention learned from what the acts did to the people they fell on. The school is where the second is formed before any landlord gets it, beside the people a score of this kind acts on, who can enroll in their own case.

The case shows what the targeting and removal cases can't: the mechanism working in full, with every channel open. It also shows the limit. The operator is private, so the custody terms apply unchanged, and the sampled re-adjudication, the discount curve and the auto-proceed partition are all measurable by an auditor; what is still unmeasured is whether the system formed on the notice channel comes to weigh the claims of the people who contested, which is the formed-attention test's question and not something a worked case can settle.
